\section{積分と漸化式}\label{節:積分漸化式}

積分に自然数$n$が含まれているとき,\,$n=0,1,2,\ldots$の値を1つずつ計算するのは手間である.例えば$\int_0^1x^n\e^{-x}\dd x$を直接求めるには,部分積分を$n$回繰り返さなければならない.そこで,\,$n$を1つ小さくした積分との関係,つまり漸化式を作ることを目標とする.

手がかりは部分積分である.\,$x^n$を微分すると$nx^{n-1}$となり,指数が1つ下がる.相手の関数$g(x)$の原始関数を$G(x)$とすると,\,$n\ge1$で
\[
 \int_a^b x^ng(x)\dd x
 =\left[x^nG(x)\right]_a^b
 -n\int_a^b x^{n-1}G(x)\dd x
\]
となる.\,$g(x)=\e^{-x}$のように,\,$G(x)$が$g(x)$の定数倍になる関数なら,右辺の積分は左辺の$n$を$n-1$に変えた積分の定数倍である.こうして隣り合う$n$の積分が漸化式で結ばれる.

\begin{bookpoint}
\textbf{POINT}\par
$x^n$を微分する側に選んで部分積分し,\,$n$番目の積分を$n-1$番目の積分で表す.最初の値は直接計算する.
\end{bookpoint}

\begin{example}{積分と漸化式}
0以上の整数$n$に対して
\[
 A_n=\int_0^1 x^n\e^{-x}\dd x
\]
とおく.\,$A_n$の漸化式を求め,\,$A_2$を計算せよ.また$A_n\to0$を示せ.
\end{example}

\paragraphbf{方針}
$x^n$を微分し,\,$\e^{-x}$を積分する部分積分で$A_n$を$A_{n-1}$で表す.極限は漸化式を使わず,\,$\e^{-x}\le1$で積分を上から抑える.

\paragraphbf{解答}
$A_n=nA_{n-1}-\e^{-1}$\,$(n\ge1)$,\,$A_0=1-\e^{-1}$,\,$A_2=2-5\e^{-1}$.\,$A_n\to0$も成り立つ.

\paragraphbf{解法}
$n\ge1$のとき,\,$x^n$を微分し,\,$\e^{-x}$を積分する部分積分を行うと
\begin{align*}
 A_n
 &=\left[-x^n\e^{-x}\right]_0^1
   +n\int_0^1x^{n-1}\e^{-x}\dd x\\
 &=-\e^{-1}+nA_{n-1}
\end{align*}
となる.\,$x=0$の端点の項が$0$になるのは$n\ge1$だからである.最初の値は
\[
 A_0=\int_0^1\e^{-x}\dd x=1-\e^{-1}
\]
と直接求める.漸化式より
\begin{align*}
 A_1&=1\cdot A_0-\e^{-1}=1-2\e^{-1},\\
 A_2&=2A_1-\e^{-1}=2-5\e^{-1}
\end{align*}
を得る.

極限を調べるには,漸化式より積分の評価が短い.\,$[0,1]$では$0<\e^{-x}\le1$なので
\[
 0<A_n\le\int_0^1x^n\dd x=\frac1{n+1}
\]
である.右辺は$0$へ近づくので,はさみうちの原理より$A_n\to0$である.

\paragraphbf{漸化式を解く}
得られた漸化式
\[
 A_n=nA_{n-1}-\e^{-1}
\]
を一般項まで解いてみよう.\,$A_{n-1}$に掛かる倍率$n$が添字とともに変わるので,特性方程式型のように定数を引いて等比型にすることはできない.そこで,倍率$n,n-1,\ldots,1$をまとめた$n!$で両辺を割ってやる.
\[
 \frac{A_n}{n!}=\frac{A_{n-1}}{(n-1)!}-\frac{\e^{-1}}{n!}
\]
ここで$B_n=\frac{A_n}{n!}$と置くと
\[
 B_n=B_{n-1}-\frac{\e^{-1}}{n!}
\]
となり,隣り合う項の差が$n$の式で与えられる階差型に帰着する.\,$k=1,\ldots,n$について足せば途中の項が消え,
\[
 B_n=B_0-\e^{-1}\sum_{k=1}^{n}\frac1{k!}
\]
を得る.\,$B_0=A_0=1-\e^{-1}$であり,\,$0!=1$なので
\[
 A_n=n!\left(1-\e^{-1}\sum_{k=0}^{n}\frac1{k!}\right)
\]
となる.ただし,この形から$n\to\infty$の極限を求めるには,括弧の中が$0$へ近づく速さと$n!$が大きくなる速さを比べなければならない.先ほどの積分による評価は,その比較をせずに直接$A_n\to0$を示している.

\paragraphbf{漸化式で値を求める実験}
漸化式があれば,\,$A_0$から順に値を計算できる.そこで$\e^{-1}=0.367879\ldots$を小数第4位までに丸めた$0.3679$を使い,\,$A_0=1-0.3679$から計算してみよう.比べるのは,積分の値を小数第4位まで求めたものである.
\begin{center}
\begin{tabular}{crr}
\toprule
$n$ & 積分の値 & 漸化式で求めた値\\
\midrule
$0$ & $0.6321$ & $0.6321$\\
$1$ & $0.2642$ & $0.2642$\\
$2$ & $0.1606$ & $0.1605$\\
$3$ & $0.1139$ & $0.1136$\\
$4$ & $0.0878$ & $0.0865$\\
$5$ & $0.0713$ & $0.0646$\\
$6$ & $0.0599$ & $0.0197$\\
$7$ & $0.0517$ & $-0.2300$\\
$8$ & $0.0454$ & $-2.2079$\\
\bottomrule
\end{tabular}
\end{center}
被積分関数は正なので$A_n>0$のはずだが,漸化式で求めた値は$n=7$で負になってしまった.

原因は丸めたことによるずれ$\delta=0.3679-\e^{-1}$である.漸化式で求めた値を$a_n$とし,\,$D_n=a_n-A_n$と置く.2つの漸化式
\[
 \begin{array}{rrl}
 &a_n={}&na_{n-1}-0.3679\\
 -\,)&A_n={}&nA_{n-1}-\e^{-1}\\ \hline
 &D_n={}&nD_{n-1}-\delta
 \end{array}
\]
の引き算から,\,$D_n$も同じ形の漸化式を満たす.\,$D_0=-\delta$なので,先ほどと同じく$n!$で割って足せば
\[
 D_n=-\delta\,n!\sum_{k=0}^{n}\frac1{k!}
\]
となる.\,$\delta$は$0.00002$程度にすぎないが,\,$n!$倍されて現れたのだ.\,$n=7$では$7!=5040$であり,ずれは約$0.28$になる.漸化式は値を順に求めるのに役立ち,積分表示は符号や大きさを調べるのに役立つ.同じ数列でも,問いに応じて使いやすい表し方が違う.

\paragraphbf{ネイピア数の級数}
一般項と積分による評価を組み合わせると,ネイピア数$\e$を和で表せる.一般項の両辺を$n!$で割って$\e$を掛けると
\[
 \frac{\e A_n}{n!}=\e-\sum_{k=0}^{n}\frac1{k!}
\]
となる.\,$0<A_n\le\frac1{n+1}$より
\[
 0<\e-\sum_{k=0}^{n}\frac1{k!}\le\frac{\e}{(n+1)!}
\]
である.右辺は$0$へ近づくので
\[
 \lim_{n\to\infty}\sum_{k=0}^{n}\frac1{k!}=\e
\]
を得る.例えば$n=5$なら和は$2.71666\ldots$であり,\,$\e=2.71828\ldots$との差は$\frac{\e}{6!}<0.004$以内に収まっている.
